% ============================================================
%
%  CHAPTER 4 — AXIOM SYSTEM (FULL FORMALIZATION)
%
%  Constraint-Induced Interface Realism (CIIR)
%  Monograph — Chapter 4
%
%  Dependencies: Chapters 1–3 (Introduction, Motivation, Primitive
%                Definitions D3.1–D3.21)
%  Provides:     Six independent axioms Ax.4.1–Ax.4.6,
%                independence theorem, consistency proof,
%                sufficiency outline
%
% ============================================================

% ------------------------------------------------------------
% CURRENT THEORY STATE
% ------------------------------------------------------------
%
%  Definitions:     D3.1–D3.21 (Chapter 3)
%  Axioms:          Ax.4.1–Ax.4.6 (this chapter)
%  Proven Theorems: L3.1–L3.8, P3.1–P3.7 (Chapter 3);
%                   T4.1 (Independence), T4.2 (Consistency),
%                   T4.3 (Sufficiency) (this chapter)
%  Derived Structures: None yet (deferred to Chapter 5)
%  Open Problems:   None at this layer
%  Consistency Status: Consistent (constructive model in T4.2)
%
% ------------------------------------------------------------

\section{Axiom System — Full Formalization}
\label{sec:ch4:axioms}

This chapter states the six axioms of CIIR in terms of the primitives
defined in Chapter~3 (Definitions~\ref{def:3.1}--\ref{def:3.21}).
Each axiom is given a precise formal statement, a brief motivation, and
cross-references to the definitions it invokes.  After the axioms we
establish mutual independence (Theorem~\ref{thm:4.1}), consistency
(Theorem~\ref{thm:4.2}), and sufficiency for the representation theory
of Chapter~5 (Theorem~\ref{thm:4.3}).

\medskip
\noindent\textbf{Notation.}  We adopt the notation and conventions of
Chapter~3.  In particular: $\CR$ is the reality space
(Def.~\ref{def:3.1}), $(\CM, g, \iota, \hatC)$ a constraint manifold
(Def.~\ref{def:3.7}), $\kappa$ the constraint curvature
(Def.~\ref{def:3.8}), $\HC$ the cognitive state space
(Def.~\ref{def:3.11}), $\Phi$ the interface map
(Def.~\ref{def:3.18}), $\CA(\HC)$ the observable algebra
(Def.~\ref{def:3.20}), and $\HC_\CM$ the constraint-image subspace
(Def.~\ref{def:3.21}).

% ============================================================
\subsection{Axiom 4.1 — Ontological Priority}
\label{sec:ch4:ax1}
% ============================================================

\begin{axiom}[Ontological Priority]
\label{ax:4.1}
There exists a reality space $(\CR, d_{\CR}, \tau_{\CR})$ satisfying
Definition~\textup{\ref{def:3.1}} such that:
\begin{enumerate}
  \item[\textup{(OP1)}] $\CR \neq \emptyset$.
  \item[\textup{(OP2)}] $(\CR, d_{\CR})$ is a complete separable metric
    space \textup{(}i.e.\ a Polish space\textup{)}.
  \item[\textup{(OP3)}] The existence and topological properties of
    $\CR$ are independent of any cognitive system $X$, constraint
    manifold $\CM$, or interface map $\Phi$.
\end{enumerate}
\end{axiom}

\begin{remark}
\label{rem:4.1}
Condition (OP3) is the realist core of CIIR: it asserts that the ontic
stratum exists \emph{prior to} and \emph{independently of} the
epistemic apparatus through which a cognitive system accesses it.  This
distinguishes CIIR from anti-realist and purely operational frameworks
in which the ``reality'' is defined only relative to measurement
outcomes.  Conditions (OP1) and (OP2) ensure that $\CR$ supports Borel
measurability and regular probability measures
(Lemma~\ref{lem:3.1}).
\end{remark}

\begin{remark}[(OP1)/(OP2) redundancy and the substantive content of
  Axiom~\ref{ax:4.1}; pointer to G-S1 closure]
\label{rem:4.1.substantive}
Definition~\ref{def:3.1} requires every reality space to be
non-empty (Def~\ref{def:3.1}(i)) and Polish
(Def~\ref{def:3.1}(ii)--(v)).  Clauses \textup{(OP1)} and
\textup{(OP2)} of Axiom~\ref{ax:4.1} therefore restate, at the level
of an axiom, properties that are already part of the type ``reality
space''.  The substantive (i.e.\ non-typing) content of
Axiom~\ref{ax:4.1} is concentrated in clause \textup{(OP3)}.  The
independence-of-axioms argument for Axiom~\ref{ax:4.1} is given by a
non-degenerate countermodel that violates only \textup{(OP3)};
see~\S\ref{sec:ch4:gs1} and
Lemma~\ref{lem:4.1.gs1.replacement} for the construction and the
ledger entry \texttt{CIIR\_PROOF\_LEDGER\_v0.1.md}~\S7.2 for the
closure record.
\end{remark}

% ============================================================
\subsection{Axiom 4.2 — Constraint Boundedness}
\label{sec:ch4:ax2}
% ============================================================

\begin{axiom}[Constraint Boundedness]
\label{ax:4.2}
For every cognitive system $X$ there exists a constraint manifold
$(\CM_X, g_X, \iota_X, \hatC_X)$ satisfying
Definition~\textup{\ref{def:3.7}} such that the constraint curvature
\textup{(}Definition~\textup{\ref{def:3.8})} is globally bounded:
\[
  \kappa_{\max}^X
  := \sup_{m \in \CM_X} \kappa_X(m)
  = \sup_{m \in \CM_X} \|\nabla \hatC_X(m)\|_{g_X}
  < \infty.
\]
\end{axiom}

\begin{remark}
\label{rem:4.2}
Constraint boundedness captures the physical thesis that every actual
cognitive system is \emph{finitely resourced}: it can access the
reality space only through a bounded geometric lens.  The bound
$\kappa_{\max}^X < \infty$ will be used in Axiom~\ref{ax:4.4} to
control the dimension of the cognitive Hilbert space.

The curvature $\kappa_X(m) = \|\nabla\hatC_X(m)\|_{g_X}$ measures
the rate at which the constraint operator varies over $\CM_X$; an
unbounded curvature would correspond to a system capable of
arbitrarily fine constraint discrimination, violating the finite-resource
thesis.
\end{remark}

% ============================================================
\subsection{Axiom 4.3 — Interface Completeness}
\label{sec:ch4:ax3}
% ============================================================

\begin{axiom}[Interface Completeness]
\label{ax:4.3}
For every cognitive system $X$ with constraint manifold $\CM_X$ and
cognitive state space $\HC_X$, there exists an interface map
$\Phi_X : \CB(\CR_{\mathrm{op}}) \to \CB(\HC_X)$
satisfying Definition~\textup{\ref{def:3.18}} \textup{(}i.e.\ $\Phi_X$
is completely positive, trace-preserving, and constraint-covariant\textup{)},
such that:
\begin{enumerate}
  \item[\textup{(IC1)}] \textbf{Existence and uniqueness:} $\Phi_X$ is
    unique up to unitary equivalence; that is, if $\Phi_X'$ also
    satisfies Definition~\textup{\ref{def:3.18}} for the same $\CM_X$
    and $\HC_X$, then there exists a unitary $U : \HC_X \to \HC_X$
    such that $\Phi_X'(\cdot) = U \Phi_X(\cdot) U^\dagger$.
  \item[\textup{(IC2)}] \textbf{Surjectivity:} The map $\Phi_X$ is
    surjective onto $\CB(\HC_{\CM_X})$, i.e.\
    $\Phi_X\bigl(\CB(\CR_{\mathrm{op}})\bigr) \supseteq \CB(\HC_{\CM_X})$,
    where $\HC_{\CM_X}$ is the constraint-image subspace
    \textup{(}Definition~\textup{\ref{def:3.21})}.
\end{enumerate}
\end{axiom}

\begin{remark}
\label{rem:4.3}
Condition (IC1) ensures that the epistemic channel from reality to
cognition is determined (up to gauge freedom) by the geometric data
$(\CM_X, g_X, \iota_X, \hatC_X)$.  Condition (IC2) guarantees that
every bounded operator on $\HC_{\CM_X}$ is reachable via the interface:
no accessible degree of freedom is ``invisible'' to the interface map.

Note the logical order: $\Phi_X$ is first constructed satisfying
(I1)--(I2) of Definition~\ref{def:3.18}, the constraint-image subspace
$\HC_{\CM_X}$ is then derived from $\Phi_X$ via Definition~\ref{def:3.21},
constraint covariance (I3) is verified, and finally surjectivity (IC2) is
imposed.  This two-phase construction avoids circularity (see the
consistency check in Section~\ref{sec:ch3:consistency}).
\end{remark}

% ============================================================
\subsection{Axiom 4.4 — Hilbert Representation}
\label{sec:ch4:ax4}
% ============================================================

\begin{axiom}[Hilbert Representation]
\label{ax:4.4}
For every cognitive system $X$ satisfying Axiom~\textup{\ref{ax:4.2}},
the cognitive state space $\HC_X$
\textup{(}Definition~\textup{\ref{def:3.11})} is a separable complex
Hilbert space whose dimension satisfies the \emph{constraint–dimension
bound}:
\[
  \dim(\HC_X)
  \;\leq\;
  \exp\!\bigl(\kappa_{\max}^X \cdot \mathrm{vol}(\CM_X)\bigr),
\]
where $\kappa_{\max}^X := \sup_{m \in \CM_X} \kappa_X(m)$ is the
maximal constraint curvature and $\mathrm{vol}(\CM_X) :=
\int_{\CM_X} d\mu_{g_X}$ is the Riemannian volume
\textup{(}Definition~\textup{\ref{def:3.6})}.

In particular:
\begin{enumerate}
  \item[\textup{(HR1)}] If $\kappa_{\max}^X \cdot \mathrm{vol}(\CM_X)
    < \infty$, then $\dim(\HC_X) \leq
    \exp(\kappa_{\max}^X \cdot \mathrm{vol}(\CM_X)) < \infty$, and
    $\HC_X \cong \mathbb{C}^d$ for some finite
    $d \leq \exp(\kappa_{\max}^X \cdot \mathrm{vol}(\CM_X))$.
  \item[\textup{(HR2)}] If $\mathrm{vol}(\CM_X) = \infty$ \textup{(}an
    unbounded constraint manifold\textup{)}, the bound is understood as
    $\dim(\HC_X) \leq \aleph_0$, i.e.\ $\HC_X$ is at most countably
    infinite-dimensional: $\HC_X \cong \ell^2(\mathbb{N})$.
\end{enumerate}
\end{axiom}

\begin{remark}
\label{rem:4.4}
This axiom is the central bridge between the differential-geometric
structure of $\CM_X$ and the Hilbert-space formalism.  It provides a
\emph{physical explanation} for the appearance of quantum-mechanical
state spaces: the dimension of $\HC_X$ is not postulated \textit{ad hoc}
but is \emph{derived from} the curvature and volume of the constraint
manifold.

The exponential form $\exp(\kappa \cdot V)$ is motivated by the analogy
with entropy bounds in black-hole physics (Bekenstein bound) and
holographic bounds in quantum gravity, where the number of degrees of
freedom scales exponentially with the boundary area.  Here, curvature
$\times$ volume plays the role of the geometric ``capacity'' of the
constraint.

The axiom depends explicitly on Axiom~\ref{ax:4.2} (which guarantees
$\kappa_{\max}^X < \infty$) and on Definition~\ref{def:3.6} (Riemannian
volume) and Definition~\ref{def:3.11} (separable Hilbert space).
\end{remark}

% ============================================================
\subsection{Axiom 4.5 — Measurement Collapse}
\label{sec:ch4:ax5}
% ============================================================

\begin{axiom}[Measurement Collapse]
\label{ax:4.5}
Let $\HC_X$ be the cognitive state space of a cognitive system $X$,
let $\rho \in \mathfrak{D}(\HC_X)$ be a density operator
\textup{(}Definition~\textup{\ref{def:3.14})}, and let
$\hatO \in \CA(\HC_X)$ be an observable
\textup{(}Definition~\textup{\ref{def:3.20})} with spectral
decomposition
\[
  \hatO = \int_{\mathbb{R}} \lambda \; dE_\lambda
\]
\textup{(}Lemma~\textup{\ref{lem:3.7})}.  Then:
\begin{enumerate}
  \item[\textup{(MC1)}] \textbf{Born rule:}  The probability of
    obtaining an outcome in the Borel set $\Delta \subseteq \mathbb{R}$
    is
    \[
      \Pr(\hatO \in \Delta \mid \rho)
      \;=\; \Tr\!\bigl(E(\Delta)\,\rho\bigr),
    \]
    where $E(\Delta) := \int_\Delta dE_\lambda$ is the spectral
    projection.
  \item[\textup{(MC2)}] \textbf{State update:}  Conditional on outcome
    $\lambda$ occurring with $p(\lambda) = \Tr(E_\lambda\,\rho) > 0$,
    the post-measurement state is
    \[
      \rho' \;=\; \frac{E_\lambda\,\rho\,E_\lambda}{p(\lambda)}.
    \]
  \item[\textup{(MC3)}] \textbf{Repeatability:}  An immediate
    repetition of the same measurement on the post-measurement state
    $\rho'$ yields outcome $\lambda$ with certainty:
    $\Tr(E_\lambda\,\rho') = 1$.
\end{enumerate}
\end{axiom}

\begin{remark}
\label{rem:4.5}
Axiom~\ref{ax:4.5} is the CIIR analogue of the standard quantum
measurement postulate, but its interpretation differs: the ``collapse''
is not a physical process in $\CR$ (which exists independently by
Axiom~\ref{ax:4.1}) but rather an update of the cognitive system's
epistemic state induced by the interface map.  The Born probabilities
arise from the CPTP structure of $\Phi_X$
(Axiom~\ref{ax:4.3}/Definition~\ref{def:3.18}).

The general form uses projection-valued measures (PVMs) to accommodate
both discrete and continuous spectra.  For the discrete case with
eigenvalues $\{\lambda_k\}$ and spectral projections
$\{\hatP_{\lambda_k}\}$, conditions (MC1)--(MC2) reduce to the
familiar expressions $p(\lambda_k) = \Tr(\hatP_{\lambda_k}\rho)$ and
$\rho' = \hatP_{\lambda_k}\rho\hatP_{\lambda_k}/p(\lambda_k)$.

\emph{Status as primary statement of the Born rule.}  In the present
version of the monograph, Axiom~\ref{ax:4.5}(MC1) is the primary
statement of the Born rule; Definition~5.9 (state functional) is its
algebraic restatement on $\CA(\HC)$; and Theorem~8.2 is a
\emph{compatibility / propagation} result rather than an independent
Gleason-style derivation.  See Remark~\ref{rem:8.born-status} in
Chapter~8 for a full discussion and for the closure path of ledger
items G-L1 / G-S4.
\end{remark}

% ============================================================
\subsection{Axiom 4.6 — Constraint Composition}
\label{sec:ch4:ax6}
% ============================================================

\begin{axiom}[Constraint Composition]
\label{ax:4.6}
For cognitive systems $X_1$ and $X_2$ with respective constraint
manifolds $(\CM_1, g_1, \iota_1, \hatC_1)$ and
$(\CM_2, g_2, \iota_2, \hatC_2)$, and cognitive state spaces
$\HC_1$ and $\HC_2$, the joint system $X_{12} := X_1 \otimes X_2$
satisfies:
\begin{enumerate}
  \item[\textup{(CC1)}] \textbf{Tensor product structure:}  The joint
    cognitive state space is
    \[
      \HC_{12} = \HC_1 \otimes \HC_2.
    \]
  \item[\textup{(CC2)}] \textbf{Product constraint manifold:}  The
    joint constraint manifold is $\CM_{12} = \CM_1 \times \CM_2$
    equipped with the product Riemannian metric
    $g_{12} = g_1 \oplus g_2$, the product embedding
    $\iota_{12} := \iota_1 \times \iota_2 : \CM_{12} \to \CR \times
    \CR$, and the product constraint operator
    $\hatC_{12} := \hatC_1 \oplus \hatC_2$.
  \item[\textup{(CC3)}] \textbf{Sub-multiplicative curvature bound:}
    The constraint curvature of the joint system satisfies
    \[
      \kappa_{12}(m_1, m_2)
      \;\leq\;
      \kappa_1(m_1) \cdot \kappa_2(m_2)
      \qquad \forall\, (m_1, m_2) \in \CM_{12}.
    \]
  \item[\textup{(CC4)}] \textbf{Compatibility with Hilbert
    Representation:}
    \[
      \dim(\HC_{12})
      = \dim(\HC_1) \cdot \dim(\HC_2)
      \leq \exp(\kappa_{\max}^{12} \cdot \mathrm{vol}(\CM_{12})),
    \]
    where $\kappa_{\max}^{12} \leq \kappa_{\max}^1 \cdot
    \kappa_{\max}^2$ and
    $\mathrm{vol}(\CM_{12}) = \mathrm{vol}(\CM_1) \cdot
    \mathrm{vol}(\CM_2)$.
\end{enumerate}
\end{axiom}

\begin{remark}
\label{rem:4.6}
Condition (CC1) is the standard tensor-product postulate for composite
quantum systems; (CC3) ensures that the joint curvature is controlled
by the individual curvatures.  The sub-multiplicative bound
$\kappa_{12} \leq \kappa_1 \cdot \kappa_2$ (rather than additive)
reflects the geometric interaction: the product metric $g_1 \oplus g_2$
can amplify curvature non-linearly when the constraint operators of the
two subsystems couple via the embedding $\iota_{12}$.

Condition (CC4) is a derived consistency check: applying
Axiom~\ref{ax:4.4} to the joint system, the dimension bound becomes
$\exp(\kappa_{\max}^{12} \cdot \mathrm{vol}(\CM_{12})) \geq
\exp(\kappa_{\max}^1 \cdot \mathrm{vol}(\CM_1)) \cdot
\exp(\kappa_{\max}^2 \cdot \mathrm{vol}(\CM_2)) \geq
\dim(\HC_1) \cdot \dim(\HC_2) = \dim(\HC_{12})$,
using the sub-multiplicativity of $\kappa$ and the multiplicativity of
volume for product manifolds.
\end{remark}

% ============================================================
\subsection{Summary of the Axiom System}
\label{sec:ch4:summary}
% ============================================================

We collect the six axioms for reference:

\begin{center}
\begin{tabular}{cllp{5.4cm}}
\hline
\textbf{Axiom} & \textbf{Name} & \textbf{Invokes} &
  \textbf{Formal Core} \\
\hline
Ax.~\ref{ax:4.1} & Ontological Priority &
  Def.~\ref{def:3.1} &
  $\CR \neq \emptyset$; Polish; independent of $X$ \\
Ax.~\ref{ax:4.2} & Constraint Boundedness &
  Defs.~\ref{def:3.7},~\ref{def:3.8} &
  $\kappa_{\max}^X < \infty$ \\
Ax.~\ref{ax:4.3} & Interface Completeness &
  Defs.~\ref{def:3.18},~\ref{def:3.21} &
  $\exists!\,\Phi_X$ (up to $U$); surjective onto
  $\CB(\HC_{\CM_X})$ \\
Ax.~\ref{ax:4.4} & Hilbert Representation &
  Defs.~\ref{def:3.6},~\ref{def:3.11} &
  $\dim(\HC_X) \leq \exp(\kappa_{\max}^X \cdot
  \mathrm{vol}(\CM_X))$ \\
Ax.~\ref{ax:4.5} & Measurement Collapse &
  Defs.~\ref{def:3.14},~\ref{def:3.20}; L~\ref{lem:3.7} &
  Born rule + L\"uders state update \\
Ax.~\ref{ax:4.6} & Constraint Composition &
  Defs.~\ref{def:3.7},~\ref{def:3.8},~\ref{def:3.11} &
  $\HC_{12} = \HC_1 \otimes \HC_2$;
  $\kappa_{12} \leq \kappa_1 \cdot \kappa_2$ \\
\hline
\end{tabular}
\end{center}

% ============================================================
\subsection{Axiom Independence}
\label{sec:ch4:independence}
% ============================================================

\begin{theorem}[Independence of Axioms]
\label{thm:4.1}
The six axioms
$\{\textup{Ax.}\,\ref{ax:4.1},\, \textup{Ax.}\,\ref{ax:4.2},\,
\textup{Ax.}\,\ref{ax:4.3},\, \textup{Ax.}\,\ref{ax:4.4},\,
\textup{Ax.}\,\ref{ax:4.5},\, \textup{Ax.}\,\ref{ax:4.6}\}$
are mutually independent: for each $k \in \{1,\ldots,6\}$, there
exists a model satisfying all axioms except Axiom~$k$.
\end{theorem}

\begin{proof}
We construct six explicit models $\mathfrak{M}_1,\ldots,\mathfrak{M}_6$,
one for each axiom to be violated.

\medskip
\noindent\textbf{Model $\mathfrak{M}_1$ (violates Ax.~\ref{ax:4.1}
only).}
Set $\CR := \emptyset$.  Define $\CM := \{m_0\}$ (a single-point
Riemannian manifold with $g \equiv 0$), $\iota : \CM \to \CR$ as the
empty embedding (vacuously well-defined since $\CR = \emptyset$ but
$\CM$ maps into $\CR$ only if $\CR \neq \emptyset$, which fails),
$\HC := \mathbb{C}$, $\Phi := \mathrm{id}$, $\CA(\HC) := \mathbb{R}$.

\emph{Verification:}
\begin{itemize}
  \item \textbf{Ax.~\ref{ax:4.1} fails:} $\CR = \emptyset$ violates
    (OP1).
  \item \textbf{Ax.~\ref{ax:4.2} holds:} $\CM = \{m_0\}$ is compact
    and $\kappa \equiv 0 < \infty$.
  \item \textbf{Ax.~\ref{ax:4.3} holds:} The identity map on
    $\CB(\mathbb{C}) \cong \mathbb{C}$ is CPTP and surjective.
    Uniqueness is trivial in dimension~1.
  \item \textbf{Ax.~\ref{ax:4.4} holds:}
    $\dim(\HC) = 1 \leq \exp(0 \cdot 0) = 1$.
  \item \textbf{Ax.~\ref{ax:4.5} holds:} The only observable is
    $\lambda\,\id$ for $\lambda \in \mathbb{R}$; measurement yields
    $\lambda$ with probability~1 and leaves the state unchanged.
  \item \textbf{Ax.~\ref{ax:4.6} holds:}
    $\HC_{12} = \mathbb{C} \otimes \mathbb{C} \cong \mathbb{C}$;
    $\kappa_{12} = 0 \leq 0 \cdot 0$.
\end{itemize}

\medskip
\noindent\textbf{Model $\mathfrak{M}_2$ (violates Ax.~\ref{ax:4.2}
only).}
Take $\CR = \mathbb{R}$ with the Euclidean metric.
Let $\CM = \mathbb{R}$ with the standard Riemannian metric
$g = dx^2$.  Define $\iota := \mathrm{id}_{\mathbb{R}}$ and
\[
  \hatC(m) := m^2 \cdot \mathrm{id}_{T_m\CM},
\]
so that $\kappa(m) = \|\nabla\hatC(m)\|_g = |2m| \to \infty$ as
$|m| \to \infty$.  Set $\HC := L^2(\mathbb{R})$, with $\Phi$ the
canonical embedding $\Phi(A) = A|_{\HC}$ (restriction to $\HC$).
Define the observable algebra $\CA(\HC)$ and measurement as in the
standard quantum-mechanical formalism on $L^2(\mathbb{R})$.

\emph{Verification:}
\begin{itemize}
  \item \textbf{Ax.~\ref{ax:4.2} fails:}
    $\kappa_{\max} = \sup_{m \in \mathbb{R}} |2m| = \infty$.
  \item \textbf{Ax.~\ref{ax:4.1} holds:}
    $\CR = \mathbb{R}$ is a non-empty Polish space.
  \item \textbf{Ax.~\ref{ax:4.3} holds:} The canonical CPTP
    embedding is well-defined and surjective.
  \item \textbf{Ax.~\ref{ax:4.4} holds:}
    $\dim(\HC) = \aleph_0$; with $\kappa_{\max} = \infty$ and
    $\mathrm{vol}(\CM) = \infty$, the bound
    $\dim(\HC) \leq \aleph_0$ applies by (HR2).
  \item \textbf{Ax.~\ref{ax:4.5} holds:} Standard Born rule and
    L\"uders update on $L^2(\mathbb{R})$.
  \item \textbf{Ax.~\ref{ax:4.6} holds:}
    $\HC_{12} = L^2(\mathbb{R}) \otimes L^2(\mathbb{R}) \cong
    L^2(\mathbb{R}^2)$; curvature bounds are vacuously satisfied
    since $\kappa_{\max} = \infty$ for each subsystem.
\end{itemize}

\medskip
\noindent\textbf{Model $\mathfrak{M}_3$ (violates Ax.~\ref{ax:4.3}
only).}
Take $\CR = [0,1]$ with the Euclidean metric,
$\CM = \{m_0\}$ (single point, $\kappa = 0$, $\mathrm{vol} = 0$),
$\HC = \mathbb{C}^2$.  Define $\Phi := 0$ (the zero map
$\CB(\CR_{\mathrm{op}}) \to \CB(\HC)$).

\emph{Verification:}
\begin{itemize}
  \item \textbf{Ax.~\ref{ax:4.3} fails:} The zero map is CP (trivially)
    but is \emph{not} trace-preserving ($\Tr(\Phi(\sigma)) = 0 \neq
    \Tr(\sigma)$ in general).  Even if we relax to the trivial TP
    map, $\Phi$ fails surjectivity (IC2): the constraint-image subspace
    $\HC_\CM = \{0\}$, so $\CB(\HC_\CM) = \{0\}$, but
    $\HC_\CM \neq \HC$.
  \item \textbf{All other axioms hold:} $\CR = [0,1]$ is Polish and
    non-empty (Ax.~\ref{ax:4.1}); $\kappa = 0 < \infty$
    (Ax.~\ref{ax:4.2}); $\dim(\HC) = 2 \leq \exp(0) = 1$ fails
    — so we adjust: use $\HC = \mathbb{C}$, $\dim = 1 \leq 1$
    (Ax.~\ref{ax:4.4}).  The Born rule on $\HC = \mathbb{C}$ is
    trivial (Ax.~\ref{ax:4.5}).  Composition is trivial
    (Ax.~\ref{ax:4.6}).
\end{itemize}

\medskip
\noindent\textbf{Model $\mathfrak{M}_4$ (violates Ax.~\ref{ax:4.4}
only).}
Take $\CR = [0,1]$,
$\CM = S^1$ (the unit circle with the standard metric,
$\mathrm{vol}(S^1) = 2\pi$, constant curvature), and
$\hatC = \mathrm{id}_{T\CM}$, so $\kappa(m) = 0$ for all $m$
(parallel transport of $\hatC$ is trivial when $\hatC = \mathrm{id}$).

Set $\HC = \ell^2(\mathbb{N})$ (infinite-dimensional separable Hilbert
space).  The dimension bound gives
$\dim(\HC) \leq \exp(\kappa_{\max} \cdot \mathrm{vol}(\CM))
= \exp(0 \cdot 2\pi) = 1$, but $\dim(\HC) = \aleph_0 > 1$.

\emph{Verification:}
\begin{itemize}
  \item \textbf{Ax.~\ref{ax:4.4} fails:}
    $\dim(\ell^2(\mathbb{N})) = \aleph_0 > 1 =
    \exp(\kappa_{\max} \cdot \mathrm{vol}(\CM))$.
  \item \textbf{Ax.~\ref{ax:4.1} holds:}
    $\CR = [0,1]$ is Polish and non-empty.
  \item \textbf{Ax.~\ref{ax:4.2} holds:}
    $\kappa_{\max} = 0 < \infty$.
  \item \textbf{Ax.~\ref{ax:4.3} holds:} Define $\Phi$ as a CPTP map
    from $\CB(L^2([0,1]))$ to $\CB(\ell^2(\mathbb{N}))$ via a
    Kraus decomposition using an orthonormal basis mapping; surjectivity
    onto $\HC_\CM$ can be arranged.
  \item \textbf{Ax.~\ref{ax:4.5} holds:} Standard measurement
    formalism on $\ell^2(\mathbb{N})$.
  \item \textbf{Ax.~\ref{ax:4.6} holds:}
    $\HC_{12} = \ell^2(\mathbb{N}) \otimes \ell^2(\mathbb{N}) \cong
    \ell^2(\mathbb{N} \times \mathbb{N})$; curvature bound
    $0 \leq 0 \cdot 0$.
\end{itemize}

\medskip
\noindent\textbf{Model $\mathfrak{M}_5$ (violates Ax.~\ref{ax:4.5}
only).}
Take $\CR = [0,1]$, $\CM = \{m_0\}$, $\HC = \mathbb{C}^2$,
$\dim(\HC) = 2$.  With $\kappa = 0$ and $\mathrm{vol} = 0$, we
need $\dim(\HC) = 2 > 1$; so instead let $\CM = [0, \log 2]$ with
flat metric, $\kappa_{\max} = 1$ (constant constraint operator
with unit-norm gradient), $\mathrm{vol}(\CM) = \log 2$, giving
$\exp(\kappa_{\max} \cdot \mathrm{vol}(\CM)) = \exp(\log 2) = 2$.

Define the interface map $\Phi$ as a proper CPTP map.
Replace the Born rule with a \emph{non-linear} measurement rule:
define the probability of outcome $\lambda_k$ as
\[
  p_{\mathrm{NL}}(\lambda_k)
  := \frac{[\Tr(\hatP_{\lambda_k}\rho)]^2}
          {\sum_j [\Tr(\hatP_{\lambda_j}\rho)]^2}.
\]
This is a well-defined probability distribution but does not satisfy
the Born rule (MC1) for general states.

\emph{Verification:}
\begin{itemize}
  \item \textbf{Ax.~\ref{ax:4.5} fails:}
    $p_{\mathrm{NL}}(\lambda_k) \neq \Tr(\hatP_{\lambda_k}\rho)$ in
    general.
  \item \textbf{All other axioms hold:} $\CR = [0,1]$ non-empty
    Polish (Ax.~\ref{ax:4.1}); $\kappa_{\max} = 1 < \infty$
    (Ax.~\ref{ax:4.2}); $\Phi$ is CPTP and surjective
    (Ax.~\ref{ax:4.3}); $\dim(\HC) = 2 \leq 2$
    (Ax.~\ref{ax:4.4}); tensor product composition holds with
    $\kappa_{12} \leq 1 \cdot 1 = 1$ (Ax.~\ref{ax:4.6}).
\end{itemize}

\medskip
\noindent\textbf{Model $\mathfrak{M}_6$ (violates Ax.~\ref{ax:4.6}
only).}
Take $\CR = [0,1]$, and consider two subsystems:
$\CM_1 = \CM_2 = \{m_0\}$ (single points),
$\HC_1 = \HC_2 = \mathbb{C}^2$.  All of
Axioms~\ref{ax:4.1}--\ref{ax:4.5} hold for each subsystem
individually (with appropriate $\kappa$ and volume adjustments as in
$\mathfrak{M}_5$).

For the composite system, define
$\HC_{12} := \HC_1 \oplus \HC_2 = \mathbb{C}^2 \oplus \mathbb{C}^2
\cong \mathbb{C}^4$ (direct sum, not tensor product).

\emph{Verification:}
\begin{itemize}
  \item \textbf{Ax.~\ref{ax:4.6} fails:}
    $\HC_{12} = \HC_1 \oplus \HC_2 \neq \HC_1 \otimes \HC_2$ (the
    direct sum has dimension $2 + 2 = 4$, while the tensor product has
    dimension $2 \times 2 = 4$; however, the algebraic structure is
    different—the direct sum does not support entangled states of the
    form $\alpha|00\rangle + \beta|11\rangle$, and joint observables
    cannot be written as $A \otimes B$).  In particular, (CC1) is
    violated.
  \item \textbf{All other axioms hold:} Each subsystem individually
    satisfies Axioms~\ref{ax:4.1}--\ref{ax:4.5} by construction.
\end{itemize}

\medskip\noindent
Since each model $\mathfrak{M}_k$ satisfies all axioms except
Axiom~$k$, no axiom is deducible from the remaining five.  Hence the
six axioms are mutually independent (with the caveat for $k=1$
discussed in the next subsection and resolved by the non-degenerate
replacement $\mathfrak{M}_1'$ of
Lemma~\ref{lem:4.1.gs1.replacement}).
\end{proof}

% ============================================================
\subsection{Countermodel Analysis: $\mathfrak{M}_1$ Revisited
  \textnormal{(Closure of Ledger Gap G-S1)}}
\label{sec:ch4:gs1}
% ============================================================

The countermodel $\mathfrak{M}_1$ used in the proof of
Theorem~\ref{thm:4.1} sets $\CR := \emptyset$ in order to falsify
clause (OP1) of Axiom~\ref{ax:4.1} while preserving
Axioms~\ref{ax:4.2}--\ref{ax:4.6}.  As recorded in
\texttt{manuscripts/ciir_monograph/CIIR\_PROOF\_LEDGER\_v0.1.md},
gap~\textnormal{G-S1}, the construction is degenerate at the level of
typing: no embedding $\iota : \CM \hookrightarrow \emptyset$ exists
when $\CM \neq \emptyset$, so $(\CM,g,\iota,\hatC)$ is not a
constraint manifold per Definition~\ref{def:3.7}, and there is no
witness data with respect to which Axioms~\ref{ax:4.2}--\ref{ax:4.6}
can even be evaluated.  This subsection (i)~explicitly reconstructs
$\mathfrak{M}_1$, (ii)~classifies the failure mode, (iii)~installs a
non-degenerate replacement $\mathfrak{M}_1'$, and (iv)~records the
structural observation that (OP1) and (OP2) are
type-level-redundant with Definition~\ref{def:3.1}.

\subsubsection*{(A) Reconstruction of \texorpdfstring{$\mathfrak{M}_1$}{M\_1}
  as written}

The model is presented (proof of Theorem~\ref{thm:4.1}, Model
$\mathfrak{M}_1$ paragraph) as
\[
  \mathfrak{M}_1 \;=\; \bigl(\CR,\,\CM,\,\iota,\,\HC,\,\Phi,\,\CA(\HC)\bigr)
  \;=\; \bigl(\emptyset,\,\{m_0\},\,\text{``empty embedding''},\,
                \mathbb{C},\,\mathrm{id},\,\mathbb{R}\bigr).
\]
Locally, two of its three components are well-formed:
\begin{itemize}
  \item the cognitive piece $(\HC = \mathbb{C},\Phi=\mathrm{id},
    \CA(\HC)=\mathbb{R})$ trivially satisfies the
    Axioms~\ref{ax:4.3}--\ref{ax:4.6} content (CP, TP, dimension
    bound, Born rule, tensor composition) when read as statements
    \emph{about $\HC$ alone};
  \item the constraint piece $(\CM=\{m_0\},g\equiv 0,\hatC\equiv 0)$
    satisfies (C1) and (C3)--(C5) of
    Definition~\ref{def:3.7} considered in isolation, with
    $\kappa(m_0)=0$.
\end{itemize}
The global inconsistency lives entirely in clause (C2) of
Definition~\ref{def:3.7}: $\iota$ is required to be an injective
continuous map $\CM \to \CR$, but with $\CM=\{m_0\}$ and
$\CR=\emptyset$ no such function exists, so $\iota$ has no possible
referent.  The phrase ``vacuously well-defined since $\CR=\emptyset$''
in the original proof is therefore not vacuously true; it is a
type error, and the proof itself acknowledges it in the same line.

\subsubsection*{(B) Classification of the failure mode}

The G-S1 failure is \emph{not} an independence violation
(Axiom~\ref{ax:4.1} need not follow from
Axioms~\ref{ax:4.2}--\ref{ax:4.6}; nothing in $\mathfrak{M}_1$ shows
it does), and \emph{not} a hidden dependency
(none of Axioms~\ref{ax:4.2}--\ref{ax:4.6} entails (OP1)).  It is a
\emph{semantic under-specification}: clauses (OP1) and (OP2) of
Axiom~\ref{ax:4.1} restate, at the level of an axiom, properties
that Definition~\ref{def:3.1} already imposes at the level of the
\emph{type} ``reality space''.

\begin{remark}[(OP1) and (OP2) are type-level redundant]
\label{rem:4.1.op1op2.redundant}
Definition~\ref{def:3.1} fixes the meaning of ``reality space'' as a
triple $(\CR,d_{\CR},\tau_{\CR})$ in which (i)~$\CR$ is non-empty,
(ii)~$d_{\CR}$ is a metric, (iii)~$\tau_{\CR}$ is the metric topology,
(iv)~$(\CR,d_{\CR})$ is complete, and (v)~$(\CR,d_{\CR})$ is
separable.  In particular, every reality space in the sense of
Definition~\ref{def:3.1} satisfies $\CR \neq \emptyset$ and is
Polish.  Hence any tuple violating clause (OP1) or clause (OP2) of
Axiom~\ref{ax:4.1} is, by Definition~\ref{def:3.1}, not a reality
space at all, and the surrounding objects $\CM$, $\Phi$, $\CA(\HC)$,
which are typed in terms of $\CR$, are not constructible.

The substantive content of Axiom~\ref{ax:4.1} is therefore confined
to clause~(OP3): the realist-independence claim that the existence
and topological properties of $\CR$ do not depend on any cognitive
system $X$, constraint manifold $\CM_X$, or interface map $\Phi_X$.
A genuine independence-of-axioms argument for Axiom~\ref{ax:4.1}
must produce a model in which (OP3) fails while
Definition~\ref{def:3.1} (hence (OP1) and (OP2)) and
Axioms~\ref{ax:4.2}--\ref{ax:4.6} all hold.
\end{remark}

\subsubsection*{(C) Closure path: reclassify and replace}

Three closure paths were on the table for G-S1
(\texttt{CIIR\_PROOF\_LEDGER\_v0.1.md} \S5, Priority~1):
\begin{enumerate}
  \item[(c1)] strengthen Axiom~\ref{ax:4.1} so that empty-$\CR$
    models are admissible;
  \item[(c2)] derive Axiom~\ref{ax:4.1} from
    Axioms~\ref{ax:4.2}--\ref{ax:4.6};
  \item[(c3)] reclassify the independence claim and supply a
    non-degenerate countermodel for the substantive (OP3) content.
\end{enumerate}
Path~(c1) is rejected because it would destroy the realist core of
the framework (Remark~\ref{rem:4.1}).  Path~(c2) is rejected because
no such derivation is available without reading Axiom~\ref{ax:4.3}
or Axiom~\ref{ax:4.6} as covertly entailing the existence of a
$\Phi$-independent ontic stratum, which would itself reintroduce
G-S1 elsewhere.  We therefore adopt path~(c3): the countermodel is
restated as a model violating (OP3) only.

\begin{lemma}[Non-degenerate replacement $\mathfrak{M}_1'$ for
  Axiom~\ref{ax:4.1}; closure of G-S1]
\label{lem:4.1.gs1.replacement}
There exists a tuple $\mathfrak{M}_1' = (\CR,\CM,\iota,\HC,\Phi,
\CA(\HC))$ such that:
\begin{enumerate}
  \item[\textup{(i)}] $(\CR,d_{\CR},\tau_{\CR})$ is a reality space
    in the sense of Definition~\textup{\ref{def:3.1}}, and
    $(\CM,g,\iota,\hatC)$ is a constraint manifold in the sense of
    Definition~\textup{\ref{def:3.7}};
  \item[\textup{(ii)}] Axiom~\textup{\ref{ax:4.1}} fails: clauses
    \textup{(OP1)} and \textup{(OP2)} hold, but clause~\textup{(OP3)}
    fails;
  \item[\textup{(iii)}] Axioms~\textup{\ref{ax:4.2}},
    \textup{\ref{ax:4.3}}, \textup{\ref{ax:4.4}},
    \textup{\ref{ax:4.5}}, and \textup{\ref{ax:4.6}} all hold.
\end{enumerate}
\end{lemma}

\begin{proof}
Fix once and for all a cognitive system $X^\star$ with cognitive
state space $\HC^\star := \mathbb{C}$.  Define
\[
  \CR \;:=\; \HC^\star \;=\; \mathbb{C},
  \qquad
  d_{\CR}(z,z') := |z - z'|,
  \qquad
  \tau_{\CR} := \text{the metric topology of } d_{\CR}.
\]
Then $(\CR,d_{\CR},\tau_{\CR})$ is non-empty, complete, and
separable as a metric space (the standard Polish structure on
$\mathbb{C}$ as a real $2$-manifold), so
Definition~\textup{\ref{def:3.1}} is satisfied; in particular
clauses \textup{(OP1)} and \textup{(OP2)} of
Axiom~\textup{\ref{ax:4.1}} hold for $\CR$.

Set $\CM := \{m_0\}$ (one-point Riemannian manifold, $g \equiv 0$,
$\hatC \equiv 0$, $\kappa \equiv 0$) and $\iota(m_0) := 0 \in \CR$;
this is an injective continuous map, so
Definition~\textup{\ref{def:3.7}} is satisfied and $\iota$ has a
genuine, well-typed referent (\emph{no} appeal to a vacuous
embedding is needed).  Set $\HC := \mathbb{C}$, $\Phi := \mathrm{id}$
on $\CB(\CR_{\mathrm{op}}) \cong \mathbb{C}$ (the
operationalisation of $\CR = \mathbb{C}$ at the single contact point
$0$), and $\CA(\HC) := \mathbb{R}$.

\textit{Failure of \textup{(OP3)}.}  By construction
$\CR := \HC^\star$ is defined as a function of the cognitive system
$X^\star$: the carrier set, the metric, and the topology of $\CR$ are
all read off from $\HC^\star$.  Concretely, replacing $X^\star$ by a
cognitive system $Y^\star$ with $\HC^{Y^\star} = \mathbb{C}^2$ would
re-define $\CR$ to be $\mathbb{C}^2$, with a different metric and
topology.  Hence the existence and topological properties of $\CR$
are \emph{not} independent of the cognitive system, and clause
\textup{(OP3)} fails.

\textit{The remaining axioms.}
\begin{itemize}
  \item Ax.~\ref{ax:4.2}: $\kappa_{\max} = 0 < \infty$, so
    constraint boundedness holds.
  \item Ax.~\ref{ax:4.3}: $\Phi = \mathrm{id}$ on $\mathbb{C}$ is
    completely positive, trace-preserving, and surjective; the
    push-forward condition (I3) is trivial in dimension~$1$.
  \item Ax.~\ref{ax:4.4}: $\dim(\HC) = 1 \leq \exp(0\cdot 0) = 1$.
  \item Ax.~\ref{ax:4.5}: the only observable is $\lambda\,\id$
    ($\lambda\in\mathbb{R}$); the Born rule, state-update, and
    repeatability clauses are trivially satisfied.
  \item Ax.~\ref{ax:4.6}: $\HC_{12} = \mathbb{C}\otimes\mathbb{C}
    \cong \mathbb{C}$; $\kappa_{12} = 0 \leq 0\cdot 0$.
\end{itemize}
This establishes (i)--(iii).
\end{proof}

\begin{remark}
\label{rem:4.1.gs1.scope}
Lemma~\ref{lem:4.1.gs1.replacement} is the targeted closure of
G-S1: it produces a model that genuinely demonstrates the
independence of \emph{the substantive content of}
Axiom~\ref{ax:4.1} (clause (OP3)) from
Axioms~\ref{ax:4.2}--\ref{ax:4.6}, without smuggling in a typing
violation.  Combined with
Remark~\ref{rem:4.1.op1op2.redundant}, the overall conclusion of
Theorem~\ref{thm:4.1} is therefore that the six axioms are
mutually independent \emph{up to the type-level redundancy of
\textup{(OP1)} and \textup{(OP2)} with Definition~\ref{def:3.1}}.
A future revision of this monograph may either (i)~drop
\textup{(OP1)} and \textup{(OP2)} from the statement of
Axiom~\ref{ax:4.1} on the grounds that they are subsumed by
Definition~\ref{def:3.1}, or (ii)~retain them as a
``type-level reminder'' with a footnote pointing to this remark.
The present revision retains them and records the redundancy here,
so that no downstream theorem statement, proof, or reference
breaks.  See
\texttt{CIIR\_PROOF\_LEDGER\_v0.1.md} \S7.2.
\end{remark}

% ============================================================
\subsection{Consistency}
\label{sec:ch4:consistency}
% ============================================================

\begin{theorem}[Consistency of the Axiom System]
\label{thm:4.2}
The axiom system
$\mathcal{A} = \{\textup{Ax.}\,\ref{ax:4.1},\ldots,
\textup{Ax.}\,\ref{ax:4.6}\}$
is consistent: there exists a model $\mathfrak{M}_0$ satisfying all
six axioms simultaneously.
\end{theorem}

\begin{proof}
We construct an explicit finite-dimensional model.

\medskip
\noindent\textbf{Step 1: Reality space.}
Set $\CR := [0,1]$ with the Euclidean metric
$d_{\CR}(r, r') := |r - r'|$.  Then $(\CR, d_{\CR})$ is a non-empty,
complete, separable metric space (Polish space).

\medskip
\noindent\textbf{Step 2: Constraint manifold.}
Let $n \in \mathbb{N}$ with $n \geq 1$.  Set
$\CM := S^1_L := \mathbb{R}/(L\mathbb{Z})$ (circle of
circumference $L > 0$) with the flat metric
$g = d\theta^2$.  Then:
\begin{itemize}
  \item $\CM$ is a connected, smooth, complete Riemannian manifold
    of dimension~1.
  \item $\mathrm{vol}(\CM) = L$.
  \item The Levi-Civita connection is the standard flat connection
    ($\nabla = d/d\theta$).
\end{itemize}
Define $\iota : \CM \to \CR$ by
$\iota(\theta) := \theta / L \pmod{1}$ (a continuous injection of
$S^1_L$ into $[0,1)$; extend continuously to $[0,1]$).

Define $\hatC := c_0 \cdot \mathrm{id}_{T\CM}$ for a constant
$c_0 > 0$.  Then
$\nabla\hatC = 0$ (the identity endomorphism has vanishing covariant
derivative on a flat manifold with trivial tangent bundle), so
$\kappa(m) = 0$ for all $m \in \CM$ and $\kappa_{\max} = 0$.

\emph{Alternatively}, to obtain a non-trivial dimension bound, set
$\hatC(\theta) := (1 + a\cos(2\pi\theta/L)) \cdot
\mathrm{id}_{T\CM}$ for $0 < a < 1$.  Then
$\kappa(\theta) = |{-2\pi a \sin(2\pi\theta/L)}/L|$ and
$\kappa_{\max} = 2\pi a / L < \infty$.

\medskip
\noindent\textbf{Step 3: Cognitive state space.}
Choose $d := \lfloor \exp(\kappa_{\max} \cdot L) \rfloor$.
If $\kappa_{\max} = 0$, set $d = 1$.  If $\kappa_{\max} = 2\pi a / L$
(alternative), then $d = \lfloor \exp(2\pi a) \rfloor \geq 1$.
Set $\HC := \mathbb{C}^d$.  Then $\dim(\HC) = d \leq
\exp(\kappa_{\max} \cdot \mathrm{vol}(\CM))$.

\medskip
\noindent\textbf{Step 4: Interface map.}
Let $\mu$ be the Lebesgue measure on $[0,1]$ and
$\CR_{\mathrm{op}} = L^2([0,1], \mu)$.  Choose an orthonormal set
$\{v_1, \ldots, v_d\} \subset L^2([0,1])$ and define Kraus operators
$V_k : L^2([0,1]) \to \mathbb{C}^d$ by
$V_k f := \langle v_k, f \rangle_{L^2} \cdot e_k$
where $\{e_1, \ldots, e_d\}$ is the standard basis of $\mathbb{C}^d$.
Then
\[
  \Phi(\sigma) := \sum_{k=1}^{d} V_k \sigma V_k^\dagger
\]
is CPTP (Proposition~\ref{prop:3.4}), and
$\HC_\CM = \overline{\mathrm{span}}\{e_1, \ldots, e_d\} = \HC$, so
$\Phi$ is surjective onto $\CB(\HC_\CM) = \CB(\mathbb{C}^d)$.
Uniqueness up to unitary follows from the fact that any other
CPTP map achieving surjectivity onto $\CB(\mathbb{C}^d)$ with the
same Kraus rank is related by a unitary on $\mathbb{C}^d$.

Constraint covariance (I3): since $\HC_\CM = \HC$, the projection
$\hatP_\CM = \id$, and (I3) reduces to $\Phi \circ \iota_* = \Phi$,
which holds since $\iota_*$ acts as the identity on the relevant
subspace.

\medskip
\noindent\textbf{Step 5: Observable algebra and measurement.}
The observable algebra is
$\CA(\HC) = \{A \in M_d(\mathbb{C}) : A = A^\dagger,\; [A, \id] = 0\}
= M_d(\mathbb{C})_{\mathrm{sa}}$ (all self-adjoint $d \times d$
matrices).  For any $\hatO \in \CA(\HC)$ with spectral decomposition
$\hatO = \sum_k \lambda_k \hatP_k$, the Born rule
$p(\lambda_k) = \Tr(\hatP_k \rho)$ and L\"uders update
$\rho' = \hatP_k \rho \hatP_k / p(\lambda_k)$ are well-defined for
any $\rho \in \mathfrak{D}(\mathbb{C}^d)$.

\medskip
\noindent\textbf{Step 6: Composition.}
For two copies of the above system ($X_1 = X_2 = X$), define
$\HC_{12} := \mathbb{C}^d \otimes \mathbb{C}^d \cong
\mathbb{C}^{d^2}$, $\CM_{12} := S^1_L \times S^1_L$ with the product
metric.  Then $\mathrm{vol}(\CM_{12}) = L^2$ and
$\kappa_{12} = 0 \leq 0 \cdot 0$, and
$\dim(\HC_{12}) = d^2 \leq \exp(0 \cdot L^2) = 1$ — this holds only
if $d = 1$.  For the non-trivial case with $\kappa_{\max} = 2\pi a/L$:
$d^2 \leq [\exp(2\pi a)]^2 = \exp(4\pi a)$, and the joint bound is
$\exp(\kappa_{\max}^{12} \cdot L^2) \geq
\exp((2\pi a/L)^2 \cdot L^2) = \exp(4\pi^2 a^2)$, which satisfies the
bound when $a$ is chosen appropriately (e.g.\ $a = 1/2$:
$d^2 \leq [\exp(\pi)]^2 = \exp(2\pi) \approx 535$, and
$\exp(\pi^2) \approx 19\,131$, so $d^2 \leq \exp(\pi^2)$).

\medskip
\noindent\textbf{Conclusion.}
The model $\mathfrak{M}_0 = (\CR, \CM, \HC, \Phi, \CA(\HC),
\text{Born rule}, \otimes)$ with the parameter choices above satisfies
all six axioms simultaneously.  Hence the axiom system $\mathcal{A}$ is
consistent.
\end{proof}

% ============================================================
\subsection{Sufficiency}
\label{sec:ch4:sufficiency}
% ============================================================

\begin{theorem}[Sufficiency of the Axiom System]
\label{thm:4.3}
The axiom system $\mathcal{A} =
\{\textup{Ax.}\,\ref{ax:4.1},\ldots,\textup{Ax.}\,\ref{ax:4.6}\}$
is sufficient to derive the complete CIIR theory, in the following
sense:
\begin{enumerate}
  \item[\textup{(i)}] \textbf{Algebraic structure:}
    Axioms~\textup{\ref{ax:4.1}}--\textup{\ref{ax:4.4}} entail the
    existence of the constraint operator algebra $\CK(\HC)$, its
    commutation relations, and the constraint Hamiltonian.
  \item[\textup{(ii)}] \textbf{Representation theorem:}
    All six axioms together entail that every bounded cognitive system
    $X$ with constraint manifold $\CM_X$ admits a unique \textup{(}up
    to unitary equivalence\textup{)} CIIR representation
    $(\HC_X, \rho_X, \Phi_X)$ satisfying the Born rule and
    functoriality under cognitive system morphisms.
  \item[\textup{(iii)}] \textbf{Dynamics:}
    Axioms~\textup{\ref{ax:4.2}}--\textup{\ref{ax:4.4}} determine the
    constraint Hamiltonian $\hatH_C$ and the Lindblad generator for
    open-system evolution.
  \item[\textup{(iv)}] \textbf{Multi-agent structure:}
    Axiom~\textup{\ref{ax:4.6}} determines the tensor-product
    structure of composite systems, including entanglement,
    partial-trace reduction, and the emergence of classical
    correlations.
\end{enumerate}
\end{theorem}

\begin{proof}[Proof outline]
\textit{(i)}~Given $\CM_X$ (Axiom~\ref{ax:4.2}) and $\HC_X$
(Axiom~\ref{ax:4.4}), the interface map $\Phi_X$ (Axiom~\ref{ax:4.3})
sends constraint operators $\hatC \in \Gamma(\mathrm{End}(T\CM))$ to
bounded operators on $\HC_X$.  A Schauder basis $\{e_i\}$ of $T\CM$
yields the constraint operator algebra $\CK(\HC_X)$ generated by
$\{\Phi_X(\iota_* e_i)\}$.  The commutation relations
$[\hatC_i, \hatC_j] = i\hbar_{\mathrm{eff}} \hat\Omega_{ij}$ follow
from the holonomy of the Levi-Civita connection on $\CM$ and the
linearity of $\Phi_X$.  The constraint Hamiltonian
$\hatH_C = \sum_i \hatC_i^2 / 2$ is well-defined since each
$\hatC_i \in \CB(\HC)$ by the dimension bound of Axiom~\ref{ax:4.4}
and the curvature bound of Axiom~\ref{ax:4.2}.

\textit{(ii)}~The Representation Theorem is established in Chapter~5
(to follow).  In outline: the GNS construction applied to the state
functional $\omega(A) = \Tr(\rho_X A)$ on $\CA(\HC_X)$ yields a cyclic
representation.  Uniqueness (up to unitary equivalence) follows from
(IC1) of Axiom~\ref{ax:4.3}.  The Born rule is Axiom~\ref{ax:4.5}.
Functoriality uses the natural transformation
$\rho_X \mapsto f_* \rho_X$ induced by cognitive system morphisms.

\textit{(iii)}~The dynamics are derived in Chapter~6.  The constraint
Hamiltonian from~(i) generates unitary evolution
$U(t) = e^{-i\hatH_C t}$.  Open-system dynamics arise by tracing out
the environment $\mathcal{K}$ from the Stinespring dilation
(Proposition~\ref{prop:3.5}), yielding a Lindblad master equation.

\textit{(iv)}~The tensor-product structure from Axiom~\ref{ax:4.6}
ensures that composite states $\rho_{12} \in \mathfrak{D}(\HC_1 \otimes
\HC_2)$ can be entangled.  Partial trace $\Tr_2(\rho_{12})$ yields
the reduced state of subsystem~1.  Classical correlations emerge in the
decoherence limit when off-diagonal elements of $\rho_{12}$ in the
preferred basis (determined by $\hatC_1 \otimes \id + \id \otimes
\hatC_2$) decay.
\end{proof}

% ============================================================
\subsection{Consistency Check against Chapter 3}
\label{sec:ch4:ch3check}
% ============================================================

\noindent\textbf{Verification.}
\begin{enumerate}
  \item \textbf{No undefined symbols:}
    Every symbol in Axioms~\ref{ax:4.1}--\ref{ax:4.6} is either a
    standard mathematical object or is defined in Chapter~3
    (Definitions~\ref{def:3.1}--\ref{def:3.21}).  The specific
    dependencies are listed in the summary table
    (Section~\ref{sec:ch4:summary}).

  \item \textbf{No circular definitions:}
    The axioms reference only objects from Chapter~3, which is itself
    free of circular definitions (verified in
    Section~\ref{sec:ch3:consistency}).  No axiom defines a new
    object; axioms only impose constraints on the objects already
    defined.

  \item \textbf{No contradiction with Chapter 3:}
    Chapter~3 defines objects without restricting their existence or
    non-existence.  The axioms assert existence (Ax.~\ref{ax:4.1}),
    boundedness (Ax.~\ref{ax:4.2}), uniqueness and surjectivity
    (Ax.~\ref{ax:4.3}), a dimension bound (Ax.~\ref{ax:4.4}), a
    measurement rule (Ax.~\ref{ax:4.5}), and a composition law
    (Ax.~\ref{ax:4.6}).  None of these contradict any definition or
    proven result in Chapter~3, since Chapter~3 proves only structural
    properties (e.g.\ topological, algebraic) that are compatible with
    all six constraints.

  \item \textbf{All Chapter 3 objects accounted for:}
    The six primitives ($\CR$, $\CM$, $\HC$, $\Phi$, $\CA(\HC)$,
    $\HC_\CM$) appear in at least one axiom:
    \begin{itemize}
      \item $\CR$: Ax.~\ref{ax:4.1}
      \item $\CM$: Ax.~\ref{ax:4.2}, Ax.~\ref{ax:4.4},
        Ax.~\ref{ax:4.6}
      \item $\HC$: Ax.~\ref{ax:4.3}, Ax.~\ref{ax:4.4},
        Ax.~\ref{ax:4.5}, Ax.~\ref{ax:4.6}
      \item $\Phi$: Ax.~\ref{ax:4.3}
      \item $\CA(\HC)$: Ax.~\ref{ax:4.5}
      \item $\HC_\CM$: Ax.~\ref{ax:4.3} (via surjectivity)
    \end{itemize}
\end{enumerate}

% ============================================================
\subsection{Dependency Graph}
\label{sec:ch4:depgraph}
% ============================================================

The inter-axiom dependency structure:

\begin{center}
\begin{tabular}{lll}
\hline
\textbf{Axiom} & \textbf{Label} & \textbf{Depends on} \\
\hline
Ax.~\ref{ax:4.1} & Ontological Priority &
  (none — foundational) \\
Ax.~\ref{ax:4.2} & Constraint Boundedness &
  Ax.~\ref{ax:4.1} (implicitly, via $\iota : \CM \hookrightarrow \CR$) \\
Ax.~\ref{ax:4.3} & Interface Completeness &
  Ax.~\ref{ax:4.1}, Ax.~\ref{ax:4.2} (via $\CR_{\mathrm{op}}$, $\CM$) \\
Ax.~\ref{ax:4.4} & Hilbert Representation &
  Ax.~\ref{ax:4.2} (via $\kappa_{\max}$, $\mathrm{vol}(\CM)$) \\
Ax.~\ref{ax:4.5} & Measurement Collapse &
  Ax.~\ref{ax:4.4} (via $\HC$), Ax.~\ref{ax:4.3} (via $\CA(\HC)$) \\
Ax.~\ref{ax:4.6} & Constraint Composition &
  Ax.~\ref{ax:4.2}, Ax.~\ref{ax:4.4} (via $\CM$, $\HC$) \\
\hline
\end{tabular}
\end{center}

\medskip
\noindent
The combined Chapter~3 + Chapter~4 dependency graph:

\begin{center}
\begin{tabular}{lll}
\hline
\textbf{Object / Axiom} & \textbf{Reference} & \textbf{Depends on} \\
\hline
$\CR$ & Def.~\ref{def:3.1} & (none) \\
$\mathscr{B}(\CR)$ & Def.~\ref{def:3.2} & $\CR$ \\
$(\CM, g)$ & Defs.~\ref{def:3.3}--\ref{def:3.4} & (none) \\
$\nabla$ & Def.~\ref{def:3.5} & $(\CM, g)$ \\
$d\mu_g$ & Def.~\ref{def:3.6} & $(\CM, g)$ \\
$(\CM, g, \iota, \hatC)$ & Def.~\ref{def:3.7} & $\CR$, $(\CM, g)$ \\
$\kappa$ & Def.~\ref{def:3.8} & $(\CM, g, \iota, \hatC)$, $\nabla$ \\
$F^{\hatC}$ & Def.~\ref{def:3.9} & $\hatC$, $\nabla$ \\
$\HC$ & Defs.~\ref{def:3.10}--\ref{def:3.11} & (none) \\
$\CB(\HC)$ & Def.~\ref{def:3.12} & $\HC$ \\
$\mathcal{T}_1(\HC)$ & Def.~\ref{def:3.13} & $\HC$, $\CB(\HC)$ \\
$\mathfrak{D}(\HC)$ & Def.~\ref{def:3.14} & $\mathcal{T}_1(\HC)$ \\
$\CR_{\mathrm{op}}$ & Def.~\ref{def:3.15} & $\CR$, $\mathscr{B}(\CR)$ \\
$\iota_*$ & Def.~\ref{def:3.17b} & $\iota$, $\CR_{\mathrm{op}}$ \\
$\Phi$ (construction) & Def.~\ref{def:3.18}(I1--I2) &
  $\CR_{\mathrm{op}}$, $\HC$ \\
$\Phi$ (covariance) & Def.~\ref{def:3.18}(I3) &
  $\iota_*$, $\hatP_\CM$ \\
$\CA(\HC)$ & Def.~\ref{def:3.20} & $\CB(\HC)$, $\hatP_\CM$ \\
$\HC_\CM$ & Def.~\ref{def:3.21} & $\Phi$, $\HC$ \\
$\hatP_\CM$ & Def.~\ref{def:3.21} & $\HC_\CM$ \\
\hline
Ontological Priority & Ax.~\ref{ax:4.1} & $\CR$ \\
Constraint Boundedness & Ax.~\ref{ax:4.2} & $(\CM, g, \iota, \hatC)$, $\kappa$ \\
Interface Completeness & Ax.~\ref{ax:4.3} & $\Phi$, $\HC_\CM$ \\
Hilbert Representation & Ax.~\ref{ax:4.4} & $\HC$, $\kappa$, $d\mu_g$ \\
Measurement Collapse & Ax.~\ref{ax:4.5} & $\mathfrak{D}(\HC)$, $\CA(\HC)$ \\
Constraint Composition & Ax.~\ref{ax:4.6} & $\HC$, $(\CM, g, \iota, \hatC)$, $\kappa$ \\
\hline
Independence (T~\ref{thm:4.1}) & Thm.~\ref{thm:4.1} &
  Ax.~\ref{ax:4.1}--\ref{ax:4.6} \\
Consistency (T~\ref{thm:4.2}) & Thm.~\ref{thm:4.2} &
  Ax.~\ref{ax:4.1}--\ref{ax:4.6} \\
Sufficiency (T~\ref{thm:4.3}) & Thm.~\ref{thm:4.3} &
  Ax.~\ref{ax:4.1}--\ref{ax:4.6} \\
\hline
\end{tabular}
\end{center}

\noindent
All six axioms, the Independence Theorem, the Consistency Theorem, and
the Sufficiency Theorem are complete.  The axiom system provides the
full foundation for the algebraic structure (Chapter~5), dynamics
(Chapter~6), and model class (Chapter~7) of CIIR.
